\bookchapter{第四十六章：ステップアップへの道——使い手からアーキテクトへ}

\chapterintro{第37〜44章の方法論をマスターして、AIを使ったプロジェクト開発をスムーズにこなせるようになったら、次は何を目指すのでしょう？この章では、「AIツールのユーザー」から「AIシステムアーキテクト」へのステップアップの道を探ります。この段階での関心事はもはや「どうすればAIにいいコードを書かせられるか」ではなく、「どうすればAIを開発プロセス全体にうまく溶け込ませられるか」です。}

\booksection{MCPプロトコルとツール拡張}

MCP（Model Context Protocol）は、2025年に登場し、2026年にかけて徐々に普及したオープンなプロトコルで、AIモデルと外部ツールの間の通信標準を定めています。簡単に言えば、MCPはAIに「ツールを使う」力を与えます。コードを生成するだけでなく、コマンドの実行、データベースへの照会、APIの呼び出し、ファイルシステムの操作までできるようになるのです。

MCPが登場する前は、AIツールごとにツールの統合方式がバラバラで、互換性がありませんでした。MCPの登場が意味するのはこうです。MCP標準に対応したツールは1つ作れば、MCPに対応したどのAIツールからでも使える。これでツール統合のコストがぐっと下がり、AIの能力の範囲も「テキストの生成」から「アクションの実行」へと広がりました。

ステップアップを目指すユーザーにとって、MCPの意義は「AIのためにカスタムのツール能力を作れる」という点にあります。たとえば、チームに独自のデプロイシステムがあるなら、MCPツールを1つ書けば、AIがそのデプロイシステムを直接呼び出せるようになり、いちいち人がターミナルに切り替えて操作する必要がなくなります。この「AIがツールを呼び出す」能力こそ、Agentic Engineeringのインフラなのです。

\booksection{カスタムワークフローとSkills}

同じAIプログラミングの流れを何度も繰り返すようになったら、それを「カスタムワークフロー」や「Skills」として固めてしまいましょう。たとえば、チームに「まずセキュリティをチェック、次にパフォーマンス、さらに保守性をチェック」という標準のコードレビュー手順があるとします。この手順をSkillとして書いておけば、AIはコードレビューのたびに自動でこの流れに沿って進めてくれます。

カスタムワークフローの核心的価値は「一貫性」です。毎回同じ基準どおりに実行されることを保証し、会話の長さや文脈の違いでブレることがありません。これはチーム開発ではとくに大切です。メンバーみんなが同じワークフローを使えば、生まれてくるコードのスタイルや品質基準もそろっていきます。

カスタムワークフローの作り方はツールごとに違います。CursorはRulesとCustom Instructionsに対応し、Claude Codeはカスタムコマンド、ReplitはRunmeに対応しています。しかし核心の考え方は同じ。あなたのベストプラクティスを「AIが実行できる指示テンプレート」としてコード化するのです。

\booksection{マルチエージェント協働のパターン}

2026年に入って、マルチエージェント協働はホットな話題になりつつあります。核心の発想はこうです。タスクを1つのAIに任せきりにするのではなく、役割の違う複数のAIに協力させる。たとえば、1体はコードを書き、1体はテストを書き、1体はコードレビューを担当する——それぞれが持ち場に専念し、互いにチェックし合います。

マルチエージェント協働の強みは「分業」と「クロスチェック」にあります。各エージェントが自分の領域に集中できるので品質が上がり、エージェント同士の相互チェックで、自分では見落としがちな問題も見つけられます。一方、弱点もはっきりしています。調整のコストが高く、やり取りのオーバーヘッドが大きく、デバッグも複雑になるのです。

\begin{tcbinfo}{ワンポイント：ひとりでも「マルチエージェント」ができる}
個人開発者なら、複雑なマルチエージェントシステムを組む必要はありません。あなたが役割を「演じ分ける」ことで、マルチエージェント協働を再現できます。たとえば、まず1つ目の会話でAIに機能のコードを書かせます。次に新しい会話を開いて（CLAUDE.mdの文脈を渡して）、「コードレビュアーとしてこのコードをレビューして」とお願いします。最後にもう1つの会話で「レビューの指摘どおりにコードを直して」と頼みます。この「ひとり多役」方式なら、面倒なインフラなしで、マルチエージェント協働のメリットの大半を得られます。

\end{tcbinfo}

\booksection{継続学習の道}

AIプログラミングは変化の速い分野です。今日のベストプラクティスは、明日には更新が必要かもしれません。学び続けるためのヒントをいくつか挙げます。

1つ目は、二次的な解説ではなく、情報の「出どころ」を追うこと。Andrej Karpathy氏の元ツイートとCLAUDE.md、AnthropicやOpenAIの公式ドキュメント、各ツールの更新履歴——これらが一次情報です。コミュニティの記事やチュートリアルはその解説やまとめであり、価値はありますが、ズレが生じていることもあります。一次資料を読む習慣をつけましょう。

2つ目は、実践から学ぶこと。「学んでからやる」ではなく「やりながら学ぶ」です。新しい種類の問題に出会うたび（初めてのファイルアップロード、初めての外部APIの接続など）をきっかけに、そのつど関連知識を学びましょう。実践の中で学ぶほうが、場面から切り離した理論学習よりはるかに効率が高いのです。

3つ目は、ツールへの感度は保ちつつ、新しもの追いには走らないこと。1〜2か月に一度は半日ほど時間を取って、最新のツールや動向をチェックしましょう。ただし、新ツールが出るたびに乗り換えるのは禁物です。判断基準は「そのツールは、自分のボトルネックを大きく改善してくれるか」です。そうでなければ、今の組み合わせで十分です。

4つ目は、自分用のナレッジベースを作ること。複雑な問題を解いたときの対話、つまずいたときの教訓、うまくいった方法を、その都度記録しておきましょう。公開する必要はありません。これはあなた自身の「方法論ノート」です。積み重なるうちに、いちばん貴重な資産になります。

\begin{tcbpractice}{やってみよう：「つまずきノート」を始めよう}
今日、ドキュメントを1つ作りましょう（メモアプリでも.mdファイルでもOK）。そして、これまでのAIプログラミングでぜひ忘れたくない心得を3つ書いてください。「聞く前にまず考える」のようなものでも構いません。今後、つまずくたび、新しい技を覚えるたびに、1つずつ書き足していきます。半年後に見返したとき、今日の自分にきっと感謝するはずです。

\end{tcbpractice}

\begin{tcbsummary}{この章のまとめ}
\begin{itemize}
  \item MCPはAIとツールをつなぐ通信標準：AIがコマンド実行·データベース照会·API呼び出しをできるようにする
  \item カスタムワークフロー/Skillsの核心は「一貫性」：ベストプラクティスを指示テンプレートにコード化する
  \item マルチエージェント＝分業＋クロスチェック。個人開発者は「ひとり多役」で代用できる
  \item 継続学習の4つのコツ：一次情報を読む、やりながら学ぶ、新しものを追いすぎない、個人ナレッジベースを作る
\end{itemize}

\end{tcbsummary}
